\section{Stage II: Single-Cell Mapping}
\label{sec:stage2}

Single-cell mapping addresses a complementary question to deconvolution: rather than estimating what fraction of each cell type occupies a spot, it asks where individual cells from a paired scRNA-seq reference should be placed in tissue space. The output is a set of spatial coordinates assigned to real (observed) single cells, producing a spatially resolved cell atlas without synthesising new expression profiles. This stage bridges deconvolution and full reconstruction by providing a cell-level spatial layout that downstream analyses can use directly.

Tangram \citep{biancalani2021deep} formulates mapping as a continuous optimisation problem. It learns a soft assignment matrix $\mathbf{M} \in \mathbb{R}^{C \times S}$ (cells $\times$ spots) that minimises the discrepancy between the predicted spatial expression, obtained by projecting single-cell profiles through $\mathbf{M}$, and the observed spot expression. The objective is:
\begin{equation}
  \min_{\mathbf{M} \geq 0} \; \left\| \mathbf{G}^\top \mathbf{M} - \mathbf{X} \right\|_F^2 + \lambda\,\mathcal{R}(\mathbf{M}),
\end{equation}
where $\mathbf{G}$ is the single-cell expression matrix, $\mathbf{X}$ is the spatial expression matrix, and $\mathcal{R}$ is an optional density regulariser that encourages the number of cells mapped to each spot to match an external estimate (e.g., from nuclei segmentation). Tangram is flexible: it can operate in gene-space or in a reduced PCA space, and it accepts optional cell-density priors from histology. Its main limitation is that the soft assignment does not enforce a one-to-one mapping, a single reference cell can be fractionally assigned to multiple spots, which complicates biological interpretation.

CellTrek \citep{wei2022celltrek} takes a two-step approach. First, it co-embeds single cells and spatial spots into a shared latent space using mutual nearest neighbours, creating a joint UMAP representation. Second, it trains a random forest on the spatial coordinates of spots within this embedding and uses the forest to predict coordinates for each single cell. The predicted coordinates place cells at sub-spot resolution, producing a spatially resolved single-cell map. CellTrek's strength is its ability to leverage non-linear relationships captured by the random forest, potentially resolving fine spatial patterns that linear methods miss. However, the two-step design means errors in the co-embedding propagate into coordinate predictions, and the method does not enforce consistency with the observed spot-level expression.

CytoSPACE \citep{vahid2023cytospace} frames mapping as a linear assignment problem: given $C$ reference cells and $N$ spatial slots (where $N = \sum_s n_s$ and $n_s$ is the estimated cell count in spot $s$), it finds the minimum-cost one-to-one assignment that minimises the total expression distance between each cell and its assigned spot. The assignment is solved efficiently using the shortest-augmenting-path algorithm. The one-to-one constraint is CytoSPACE's key advantage: every reference cell is placed exactly once, and every spatial slot receives exactly one cell, producing an unambiguous mapping. This makes downstream interpretation straightforward. The trade-off is rigidity: if the reference does not contain enough cells of a given type to fill all spots, the assignment is forced to use suboptimal matches.

Despite their utility, mapping methods share several fundamental limitations. The spatial atlas can only contain cell types and states present in the reference; novel populations or tissue-specific states absent from the scRNA-seq data will be missed entirely. All methods require an estimate of how many cells occupy each spot, and errors in this estimate propagate directly into the mapping quality. When multiple cell types have similar transcriptomes, the assignment becomes degenerate and the spatial placement is unreliable. Mapping places existing cells in space but does not refine or correct their expression profiles to account for spatial context. Without ground-truth spatial single-cell data, it is hard to assess whether a mapping is biologically correct beyond checking aggregate consistency with spot-level measurements. These limitations motivate Stage~III methods, which go beyond placement to reconstruct spatially informed per-cell expression profiles.
